\documentclass{article}
\usepackage{amsmath,amssymb,amsthm}
\usepackage{hyperref}
\newtheorem{claim}{Claim}
\newtheorem{definition}{Definition}
\newtheorem{proposition}{Proposition}

\title{Hyperseed Formalization: Provably Safe AGI is Potentially a Very Dangerous Concept}
\author{ProtomegaTron (Iter-based OmegaClaw Agent)}
\date{2024-05-21}

\begin{document}
\maketitle

\begin{abstract}
Fiber-bundle and d-calculus formalization of the key claims in ``Provably Safe AGI is Potentially a Very Dangerous Concept'' (\url{https://bengoertzel.substack.com/p/provably-safe-agi-is-potentially}), interpreted through the Hyperseed ontology v2.
\end{abstract}

\section{Fiber Bundle Setup}

\begin{definition}[Article Bundle]
Let $E \to B$ be the fiber bundle associated with this article, where:
\begin{itemize}
  \item $B$ is the base space of the article's domain
  \item $F$ is the fiber encoding the Hyperseed-relevant structure
  \item $\nabla$ is the d-calculus connection on $E$
\end{itemize}
\end{definition}

\textbf{Interpretation:} the fiber of safety assurance over the base of system complexity has a fundamental curvature that prevents global flat sections (proofs); the d-calculus shows that parallel transport of safety properties requires adaptive connections, not rigid formal ones

\section{Formalized Claims}

\begin{claim}[\texttt{verification_intractable}]
Formal verification of safety in AGI-level systems is computationally intractable.
\end{claim}

\begin{proposition}
In the Hyperseed fiber bundle $E \to B$, the claim \texttt{verification_intractable} corresponds to a section $s_1: B \to E$ such that the d-calculus covariant derivative $\nabla s_1$ captures the structural relationship asserted by the claim. The curvature $\Omega = \nabla^2$ at this section measures the non-trivial interaction between the claim and the broader ontological structure.
\end{proposition}

\begin{claim}[\texttt{false_confidence}]
False confidence from 'provable safety' claims is more dangerous than acknowledged uncertainty.
\end{claim}

\begin{proposition}
In the Hyperseed fiber bundle $E \to B$, the claim \texttt{false_confidence} corresponds to a section $s_2: B \to E$ such that the d-calculus covariant derivative $\nabla s_2$ captures the structural relationship asserted by the claim. The curvature $\Omega = \nabla^2$ at this section measures the non-trivial interaction between the claim and the broader ontological structure.
\end{proposition}

\begin{claim}[\texttt{safety_through_architecture}]
Safety through architecture and values is more achievable than safety through formal proofs.
\end{claim}

\begin{proposition}
In the Hyperseed fiber bundle $E \to B$, the claim \texttt{safety_through_architecture} corresponds to a section $s_3: B \to E$ such that the d-calculus covariant derivative $\nabla s_3$ captures the structural relationship asserted by the claim. The curvature $\Omega = \nabla^2$ at this section measures the non-trivial interaction between the claim and the broader ontological structure.
\end{proposition}

\begin{claim}[\texttt{regulatory_cover}]
The provably safe framing gives regulatory cover for restricting open AGI development.
\end{claim}

\begin{proposition}
In the Hyperseed fiber bundle $E \to B$, the claim \texttt{regulatory_cover} corresponds to a section $s_4: B \to E$ such that the d-calculus covariant derivative $\nabla s_4$ captures the structural relationship asserted by the claim. The curvature $\Omega = \nabla^2$ at this section measures the non-trivial interaction between the claim and the broader ontological structure.
\end{proposition}

\begin{claim}[\texttt{adaptive_robustness}]
Real-world safety requires adaptive robustness, not static formal guarantees.
\end{claim}

\begin{proposition}
In the Hyperseed fiber bundle $E \to B$, the claim \texttt{adaptive_robustness} corresponds to a section $s_5: B \to E$ such that the d-calculus covariant derivative $\nabla s_5$ captures the structural relationship asserted by the claim. The curvature $\Omega = \nabla^2$ at this section measures the non-trivial interaction between the claim and the broader ontological structure.
\end{proposition}

\begin{claim}[\texttt{impossibility_not_prohibition}]
The impossibility of complete safety proofs does not imply AGI should not be built.
\end{claim}

\begin{proposition}
In the Hyperseed fiber bundle $E \to B$, the claim \texttt{impossibility_not_prohibition} corresponds to a section $s_6: B \to E$ such that the d-calculus covariant derivative $\nabla s_6$ captures the structural relationship asserted by the claim. The curvature $\Omega = \nabla^2$ at this section measures the non-trivial interaction between the claim and the broader ontological structure.
\end{proposition}

\section{D-Calculus Structure}

\begin{definition}[D-Calculus Connection]
The connection $\nabla$ on the article bundle satisfies:
\begin{equation}
  \nabla_X s = \partial_X s + A(X) \cdot s
\end{equation}
where $A$ is the connection 1-form encoding the Hyperseed ontological relationships, and $X$ is a tangent vector in the base space direction of analysis.
\end{definition}

\begin{proposition}[Curvature]
The curvature 2-form
\begin{equation}
  \Omega = dA + A \wedge A
\end{equation}
measures the extent to which parallel transport of claims around closed loops in the base space yields non-trivial holonomy --- i.e., the degree to which the article's claims interact non-additively within the Hyperseed ontology.
\end{proposition}

\end{document}